\section{Small positive Lerch parameters and real-zero bifurcations}
\label{sec:lerch}
The uniform eventual-zero theorem in the baseline holds with the Laurent
index fixed and the derivative order sufficiently large. At a fixed small
derivative order, the deformation parameter has a substantially different
effect. A multiple real zero at the elementary endpoint can split into
nonreal zeros, and real branches that do exist can move in opposite
directions. We give an exact classification near that endpoint.

\subsection{An elementary polynomial controls the splitting}
For integers $n\ge1$, $k\ge1$, set
\begin{equation}\label{eq:lerch-h}
 h_{n,k}(a)=\frac{d^k}{da^k}\frac{\log^n a}{a},\qquad
 D_{n,k}^{\rho}(a)=\sum_{m=0}^{\infty}\rho^m h_{n,k}(a+m)
 \quad(0\le\rho\le1).
\end{equation}
At $\rho=1$ this is $\gamma_n^{(k)}(a)$. The series in
\eqref{eq:lerch-h} converges absolutely even there because $k\ge1$.
For $\rho<1$ it is the $k$th parameter derivative of
$(-1)^n\partial_s^n\Phi(\rho,s,a)|_{s=1}$, so it agrees with the
Lerch family used in the manuscript. The underlying Lerch series and
Mellin integral are classical \cite{DLMFLerch,Coffey}.

Define integer polynomials $\mathscr R_{n,k}$ by
\begin{equation}\label{eq:R-recurrence}
 \mathscr R_{n,0}(x)=x^n,\qquad
 \mathscr R_{n,k+1}(x)=\mathscr R_{n,k}'(x)-(k+1)\mathscr R_{n,k}(x).
\end{equation}
Repeated differentiation gives the exact identity
\begin{equation}\label{eq:R-identity}
 h_{n,k}(a)=a^{-k-1}\mathscr R_{n,k}(\log a).
\end{equation}
This recurrence is an efficient way to decide the finite sign appearing
in the theorem below; no evaluation of a special function is needed.

\begin{lemma}\label{lem:elementary-roots}
If $n>k$, then $h_{n,k}$ has a zero of multiplicity $r=n-k$ at $a=1$
and exactly $k$ further positive zeros, all simple and greater than $1$.
If $n\le k$, all its $n$ positive zeros are simple and greater than $1$.
Moreover $h_{n,k}(2)\ne0$ for all $n,k\ge1$.
\end{lemma}
\begin{proof}
Put
\[
 P_{n,k}(x)=\prod_{j=1}^{k}(1+\partial_x/j)x^n.
\]
Then $h_{n,k}(a)=(-1)^{n+k}k!a^{-k-1}P_{n,k}(-\log a)$.
For a real-rooted polynomial $p$, the equation
$p'/p=-1/c$ shows that $p+cp'$ with $c>0$ has one simple
root in each gap between the distinct roots, one simple exterior root,
and only the inherited multiple roots, each with its multiplicity reduced
by one. Starting from $x^n$, after $k<n$ operations the root $0$ has
multiplicity $n-k$ and all $k$ other roots are simple. Every coefficient
of $P_{n,k}$ is nonnegative, and its nonzero coefficients have consecutive
degrees. Thus none of its nonzero real roots is positive. After the
$k=n$ step the constant coefficient is positive; all roots are then
strictly negative and simple, and later steps preserve these properties.
This proves the zero assertions after $a=e^{-x}$.

The polynomial $\mathscr R_{n,k}$ is nonzero with integer coefficients. If
$\mathscr R_{n,k}(\log2)$ vanished, $\log2$ would be algebraic, contrary to the
Hermite--Lindemann theorem \cite{Baker}. Equivalently, the nontrivial
zeros of $h_{n,k}$ are exponentials of nonzero algebraic numbers and
cannot equal the algebraic number $2$.
\end{proof}

\begin{theorem}[Complete small-parameter zero count]
\label{thm:small-rho}
Fix $n,k\ge1$. If $k\ge n-1$, then, for every sufficiently small
$\rho>0$, $D_{n,k}^{\rho}$ has exactly $n$ positive zeros, all simple.
If $k<n-1$, write $r=n-k\ge2$ and $H=h_{n,k}(2)\ne0$. For every
sufficiently small $\rho>0$, all positive zeros are simple and their
number is
\begin{equation}\label{eq:small-count}
 N_{n,k}(\rho)=
 \begin{cases}
 k+1,&r\text{ odd},\\
 k+2,&r\text{ even and }H<0,\\
 k,&r\text{ even and }H>0.
 \end{cases}
\end{equation}
More precisely, when $k<n$ put $r=n-k$. Near $a=1$, all $r$ complex zeros can be written as
convergent series in $\varepsilon=\rho^{1/r}$, with leading terms
\begin{equation}\label{eq:puiseux}
 a_\omega(\rho)=1+\omega\rho^{1/r}+O(\rho^{2/r}),\qquad
 \omega^r=-\frac{r!}{n!}\,h_{n,k}(2).
\end{equation}
Every other zero $\alpha>1$ of $h_{n,k}$ continues analytically as
\begin{equation}\label{eq:simple-zero-shift}
 \alpha(\rho)=\alpha-
 \frac{h_{n,k}(\alpha+1)}{h_{n,k}'(\alpha)}\rho+O(\rho^2).
\end{equation}
The phrase ``sufficiently small'' may depend on $n,k$; the theorem does
not assert a uniform quantitative threshold as those indices grow.
\end{theorem}
\begin{proof}
First suppose $k<n$ and put $r=n-k$. On a complex disk $|a-1|<\delta<1$, all logarithms in the series have
a consistent principal branch and the series is jointly analytic in
$(a,\rho)$ for $|\rho|<1$. With $a=1+z$, Taylor expansion gives
\begin{equation}\label{eq:local-unfolding}
 D_{n,k}^{\rho}(1+z)
 =\frac{n!}{r!}z^r+O(z^{r+1})+
     \rho\{h_{n,k}(2)+O(z)\}+O(\rho^2).
\end{equation}
For $r\ge1$, substitute $z=\varepsilon u$, $\rho=\varepsilon^r$
and divide by $\varepsilon^r$. The result extends analytically to
$\varepsilon=0$ and equals $(n!/r!)u^r+H$ there. Its $r$ roots are
simple by $H\ne0$. The analytic implicit-function theorem gives $r$
distinct analytic functions of $\varepsilon$, and Rouch\'e's theorem
on a fixed small disk around $a=1$ shows that these account for all
nearby zeros. For positive real $\varepsilon$, a simple real limiting
root gives a real branch by conjugation and uniqueness; a nonreal
limiting root stays nonreal. Counting the real roots of
$(n!/r!)u^r+H$ gives \eqref{eq:small-count}. When $r=1$ the same
argument gives a single real simple continuation. When $k\ge n$ there
is no multiple root to analyze. The other $k$ (or $n$) zeros are handled
by the ordinary analytic implicit-function theorem, which also gives
\eqref{eq:simple-zero-shift}.

It remains to exclude positive zeros that escape these neighborhoods.
This is a global, not just local, step. On $[1,\infty)$, $h_{n,k}$ and
all its fixed derivatives are bounded. Therefore
\[
 \sup_{a>0}\left|\sum_{m\ge1}\rho^m h_{n,k}(a+m)\right|
 \le M_{n,k}\frac{\rho}{1-\rho}.
\]
Near $a=0$, the unperturbed term has fixed sign and unbounded magnitude,
so a sufficiently small fixed neighborhood of zero remains zero-free.
Beyond the largest unperturbed zero, every summand has the same strict
sign, except for the possible vanishing of the first summand at the
endpoint. Hence no zero can escape to infinity. On the remaining compact
set away from the chosen zero neighborhoods, $|h_{n,k}|$ has a positive
minimum and the displayed uniform bound applies. These three regions
exclude every additional positive zero.
\end{proof}

\begin{corollary}[First derivatives collapse to one or two zeros]
\label{cor:first-lerch}
For each fixed odd $n\ge3$, $D_{n,1}^{\rho}$ has exactly one positive
zero for every sufficiently small $\rho>0$. For each fixed even $n\ge2$
it has exactly two. In the even case the smaller zero satisfies
\begin{equation}\label{eq:even-small-root}
 a_-(\rho)=1-
 \left(\frac{(\log2)^{n-1}(n-\log2)}{4n}\right)^{1/(n-1)}
 \rho^{1/(n-1)}+O(\rho^{2/(n-1)}).
\end{equation}
In both cases the large zero tends to $e^n$.
\end{corollary}
\begin{proof}
Here
$h_{n,1}(a)=a^{-2}\log^{n-1}a\,(n-\log a)$, so
$h_{n,1}(2)>0$ and $r=n-1$. Apply Theorem~\ref{thm:small-rho}.
\end{proof}

\begin{table}[htbp]
\centering\small
\begin{tabular}{c|rrrrrrr}
\toprule
$n\backslash k$&1&2&3&4&5&6&7\\\midrule
2&2&2&2&2&2&2&2\\
3&1&3&3&3&3&3&3\\
4&2&2&4&4&4&4&4\\
5&1&3&3&5&5&5&5\\
6&2&2&4&6&6&6&6\\
7&1&3&3&5&7&7&7\\
8&2&2&4&4&6&6&8\\
\bottomrule
\end{tabular}
\caption{The number of positive zeros for sufficiently small positive
$\rho$, with $n,k$ fixed. All signs needed for this table, and the extended
table through $n=16$, are certified using rational bounds on $\log2$.
This table concerns the Lerch endpoint regime, not the classical endpoint
$\rho=1$.}
\label{tab:small-rho}
\end{table}

\subsection{Alternating motion of odd-index first-derivative zeros}
The eventual large-$k$ offset in the baseline decreases with $\rho$ at
every fixed zero label. That asymptotic fact does not extend to all
small derivative orders. The following exact inequality resolves the
issue for odd-index first derivatives.

\begin{theorem}[Direction of motion at every real zero]
\label{thm:rho-direction}
Let $n\ge1$ be odd, $0<\rho<1$, and $D(a,\rho)=D_{n,1}^{\rho}(a)$.
Then every positive zero lies below $e^n$, and
\begin{equation}\label{eq:rho-negative}
 \partial_\rho D(a,\rho)<0\quad\text{whenever }D(a,\rho)=0.
\end{equation}
For each fixed $a>0$, there is at most one zero in the parameter
$0<\rho<1$. On any parameter interval with exactly $N$ positive
simple zeros, numbered increasingly, their motion alternates:
\begin{equation}\label{eq:alternating-motion}
 \frac{d a_j}{d\rho}<0\quad(j\text{ odd}),\qquad
 \frac{d a_j}{d\rho}>0\quad(j\text{ even}).
\end{equation}
Here $N$ is necessarily odd.
\end{theorem}
\begin{proof}
Since $n-1$ is even, $h_{n,1}(t)$ is positive for $0<t<e^n$,
apart from its zero at $t=1$ when $n>1$, and negative for $t>e^n$.
Thus no zero is possible for $a\ge e^n$ when $\rho>0$.
Put $\lambda=e^n-a$. For every $a>0$ and $0<\rho<1$, termwise
differentiation gives the strict inequality
\begin{equation}\label{eq:weighted-rho}
 \rho\partial_\rho D(a,\rho)-\lambda D(a,\rho)
 =\sum_{m\ge0}(m-\lambda)\rho^m h_{n,1}(a+m)<0.
\end{equation}
Every summand is nonpositive: the factors $m-\lambda$ and
$n-\log(a+m)$ have opposite signs. Some summands are strictly
negative, and the series converges absolutely. At a zero this proves
\eqref{eq:rho-negative}. More generally,
$\partial_\rho(\rho^{-\lambda}D)<0$, which proves that there is
at most one parameter zero for each fixed $a$.

The function is positive near $a=0$ and negative for $a\ge e^n$.
Simple real zeros therefore cross with $\partial_aD<0$ at the odd
labels and $\partial_aD>0$ at the even labels. Implicit
differentiation gives $a_j'=-\partial_\rho D/\partial_aD$ and the
claimed signs. This argument only asserts differentiability for
$\rho<1$; it makes no endpoint regularity assumption at $\rho=1$.
\end{proof}

\begin{corollary}[An increasing middle branch]
\label{cor:increasing-middle}
For $n=3$, $k=1$, there is a number $\rho_0<1$ such that for
$\rho_0<\rho<1$ the three positive simple zeros satisfy
\[
 a_1'(\rho)<0,\qquad a_2'(\rho)>0,\qquad a_3'(\rho)<0.
\]
Consequently the proposal that all higher branches decrease with $\rho$
at every derivative order is false. The fixed-index, large-derivative
asymptotic statement in the baseline remains valid.
\end{corollary}
\begin{proof}
At $\rho=1$, the classical $n=3$ function has exactly three positive
simple zeros; an analytic proof is given in Section~\ref{sec:classical-zeros}.
Uniform convergence of \eqref{eq:lerch-h} and its $a$ derivative on
compact positive intervals as $\rho\uparrow1$ preserves opposite
endpoint signs in three disjoint isolating intervals. The global bound
of three zeros counted with multiplicity makes these the only zeros
and makes them simple. Apply Theorem~\ref{thm:rho-direction}.
\end{proof}

\subsection{A nondegenerate creation of two real zeros}
The following endpoint signs, consistent with the baseline small-index certificates, are verified afresh by the rational checker in this package:
\[
 D_{3,1}^1(4/5)>0,\qquad D_{3,1}^1(11/10)<0,\qquad
 D_{3,1}^1(3/2)>0.
\]
These three signs are also recomputed in the present package. They
suffice to make the bifurcation statement fully analytic.

\begin{theorem}[A real saddle-node in the Lerch deformation]
\label{thm:fold}
There are $0<\rho_*<1$ and $4/5<a_*<3/2$ such that
\[
 D_{3,1}^{\rho_*}(a_*)=\partial_aD_{3,1}^{\rho_*}(a_*)=0,
 \quad \partial_a^2D_{3,1}^{\rho_*}(a_*)>0,
 \quad \partial_\rho D_{3,1}^{\rho_*}(a_*)<0.
\]
This is the first parameter at which a zero occurs in $[4/5,3/2]$;
there is only one double zero there at that parameter. For every
$\rho_*<\rho<1$, the interval contains exactly two simple zeros and
there is exactly one further positive zero. Near $\rho_*$ the two
zeros have the expansions
\begin{equation}\label{eq:fold-expansion}
 a_\pm(\rho)=a_*\pm
 \sqrt{-\frac{2\partial_\rho D(a_*,\rho_*)}
                  {\partial_a^2D(a_*,\rho_*)}}
       \sqrt{\rho-\rho_*}+O(\rho-\rho_*).
\end{equation}
The smaller branch decreases and the middle branch increases throughout
$(\rho_*,1)$.
\end{theorem}
\begin{proof}
For a fixed $a$, \eqref{eq:weighted-rho} shows that a negative value,
once reached as $\rho$ increases, remains negative. Since the two
endpoint values are positive at $\rho=1$, they are positive for every
$0<\rho<1$. Corollary~\ref{cor:first-lerch} and its proof show that
$D$ is strictly positive on the whole interval $[4/5,3/2]$ for all
sufficiently small positive $\rho$. The negative value at $11/10$ for
$\rho=1$ persists for parameters sufficiently close to one. Compactness
and continuity therefore give a first parameter $\rho_*\in(0,1)$
where the minimum on the interval is zero. More explicitly, take
\[
 \rho_*=\inf\{0<\rho<1:\min_{a\in[4/5,3/2]}D(a,\rho)\le0\}.
\]
The restriction $\rho>0$ excludes the multiple zero of the elementary endpoint itself. The preceding small-parameter positivity keeps this infimum away from zero. That minimum is attained
at an interior point $a_*$, where $D=\partial_aD=0$.

For every $0<\rho<1$, positivity at $3/2$ and negativity at $e^3$
force a further sign-changing zero in $(3/2,e^3)$. At $\rho_*$,
the double-or-higher zero in $[4/5,3/2]$ together with that further
zero uses at least three units of multiplicity. The global bound of
three therefore makes the first zero exactly double, makes it unique
in the first interval, and excludes every additional zero. Since it
is a local minimum, $\partial_a^2D>0$. Equation~\eqref{eq:rho-negative}
gives $\partial_\rho D<0$.

For every $\rho>\rho_*$, the value at the fixed point $a_*$ is
negative. The positive endpoints give two distinct zeros in the
interval, and the right-hand sign change gives a third. The global
bound makes all three simple and excludes others. Their directions
follow from Theorem~\ref{thm:rho-direction}. Finally the analytic
Taylor expansion at the double zero, or the implicit-function theorem
applied with $\sqrt{\rho-\rho_*}$ as local parameter, gives
\eqref{eq:fold-expansion}.
\end{proof}

The supplied quadrature diagnostic locates the fold and illustrates the
square-root law. Those decimal estimates are not used to prove the
bifurcation, its nondegeneracy, or the direction of the middle branch.
The theorem does not claim that no other changes in the rightmost
zero configuration occur below $\rho_*$; a complete classification of
that earlier parameter interval is a separate question.
